\documentclass[]{vvsu}

\vvsuyear{2025}

%%%%%%%%%%%%%%%%%%%

\usepackage{graphicx} % для изображений
\usepackage{tabularray} % для таблиц
\usepackage{siunitx} % для обозначений (процент, градус)
\usepackage{listings} % для листингов кода

% Список путей, где будут искаться изображения и файлы
\graphicspath{{images/}}

% Автор документа
\author{Т. Г. Пономаренко}

% Настройка стилей для листингов кода
\input{listing_styles.tex}

%%%%%%%%%%%%%%%%%%%

\begin{document}

% Шапка
\vvsuhead{\linespread{1}\selectfont{}МИНОБРНАУКИ РОССИИ\\
\vspace{10pt}Федеральное государственное бюджетное образовательное учреждение\\
высшего образования\\
\fontsize{13}{13}\selectfont{}<<ВЛАДИВОСТОКСКИЙ ГОСУДАРСТВЕННЫЙ УНИВЕРСИТЕТ>>\\
(ФГБОУ ВО <<ВВГУ>>)\\
\vspace{10pt}\fontsize{12}{12}\selectfont{}ИНСТИТУТ ИНФОРМАЦИОННЫХ ТЕХНОЛОГИЙ И АНАЛИЗА ДАННЫХ\\
КАФЕДРА ИНФОРМАЦИОННЫХ ТЕХНОЛОГИЙ И СИСТЕМ}

% Название отчета
\title{Отчет\\по лабораторной работе №5}
\subtitle{по дисциплине\\<<Информатика и программирование>>}

% Участники работы
\member{Студент\\ гр. БИН-25-3}{Т. Г. Пономаренко}
\member{Ассистент\\ преподавателя}{М.В. Водяницкий}

% Вывод титульника
\maketitle

% Задание
\begin{addition}{Задание}
  Выполнить задания и оформить отчет по стандартам ВВГУ.

   \textit{\textbf{Задание 1.}}  
  Написать функцию, которая конвертирует время из одной величины в другую.  
  На вход подаются:  
  \begin{vvsu_itemize}
    \item число (величина времени)
    \item исходная единица измерения
    \item единица измерения, в которую нужно перевести
  \end{vvsu_itemize}
  Функция должна вернуть конвертированное значение.

  Пример:\\
    Введите время для конвертации: 4 h m\\
    Результат: 240m

  \textit{\textbf{Задание 2.}}  
  Пользователь делает вклад в банке в размере \texttt{a} рублей сроком на \texttt{n} лет. Процент по вкладу зависит от суммы и срока.

  \textbf{Зависимость от суммы:}
  \begin{vvsu_itemize}
    \item каждые 10\,000 рублей увеличивают ставку на 0.3\%
    \item суммарное увеличение не может превышать 5\%
    \item минимальный вклад — 30\,000 рублей
  \end{vvsu_itemize}

  \textbf{Зависимость от срока:}
  \begin{vvsu_itemize}
    \item первые 3 года — 3\%
    \item от 4 до 6 лет — 5\%
    \item более 6 лет — 2\%
  \end{vvsu_itemize}

  Необходимо написать функцию, которая рассчитывает \textbf{прибыль пользователя без учета первоначально вложенной суммы}. Используется сложный процент.

  Пример:\\
    Ввод: 30000 3\\
    Прибыль: 3648.67 руб.

  \textit{\textbf{Задание 3.}}  
  Написать функцию для вывода всех простых чисел в заданном диапазоне. Нужно учитывать некорректные данные (например, начало больше конца или диапазон без простых чисел).

  На вход подаются два числа: начало и конец диапазона (включительно).  
  На выходе — список всех простых чисел или сообщение об ошибке.

  Пример:\\
    Начало диапазона: 1\\
    Конец диапазона: 10\\
    2 3 5 7

  Пример:\\
    Начало диапазона: 0\\
    Конец диапазона: 1\\
    Error!

  \textit{\textbf{Задание 4.}}  
  Реализовать функцию сложения двух матриц.  
  При сложении двух матриц получается новая матрица того же размера, где каждый элемент — это сумма элементов с тем же индексом из двух исходных матриц.

  \textbf{Ограничения:}
  \begin{vvsu_itemize}
    \item складывать можно только матрицы одинакового размера
    \item размер матрицы должен быть строго больше 2 (например, 3×3, 4×4 и т.д.)
    \item при нарушении условий нужно вывести сообщение об ошибке
  \end{vvsu_itemize}

  На вход подаются:
  \begin{vvsu_list}
    \item размер матрицы \texttt{n} (для квадратной матрицы \texttt{n × n})
    \item элементы первой матрицы (по строкам, через пробел)
    \item элементы второй матрицы (в таком же формате)
  \end{vvsu_list}

  Результат — новая матрица (в том же формате), либо сообщение об ошибке.

  \textit{\textbf{Задание 5.}}  
  Написать функцию, которая определяет, является ли строка палиндромом.  
  Палиндром — это строка, которая читается одинаково слева направо и справа налево (без учета пробелов, регистра и знаков препинания).

  На вход подается строка.  
  На выходе:
  \begin{vvsu_itemize}
    \item \texttt{Да}, если это палиндром
    \item \texttt{Нет}, если это не палиндром
  \end{vvsu_itemize}

  Пример:\\
    А роза упала на лапу Азора\\
    Да

  Пример:\\
    Алфавитный порядок\\
    Нет
\end{addition}

% Содержание
\toc

% Глава - Выполнение работы
\section{Выполнение работы}

% Подглава - Задание 1
\subsection{Задание 1}

Программа ищет все числа 3 в списке и заменяет их на число 30. Она работает через цикл по индексам списка: на каждом шаге проверяется, равен ли текущий элемент числу 3. Если условие выполняется, значение по этому индексу меняется на 30. В результате получается модифицированный список, где все тройки заменены на новое значение.На рисунке \ref{fig:code_task_1} представлен код программы.

\begin{vvsu_figure}{Листинг программы для задания 1}{fig:code_task_1}
  \begin{minipage}{.75\textwidth}
    \lstinputlisting[language=Python,basicstyle=\fontsize{10}{10}\linespread{1}\selectfont\ttfamily]{code/task1.py}
  \end{minipage}
\end{vvsu_figure}

% Подглава - Задание 2
\subsection{Задание 2}
Программа изменяет исходный список, возводя каждый его элемент в квадрат. Она проходит по всем индексам списка с помощью цикла, вычисляет квадрат текущего элемента и записывает результат на его место. В итоге выводится новый список, состоящий из квадратов исходных чисел. На рисунке \ref{fig:code_task_2} представлен код программы.

\begin{vvsu_figure}{Листинг программы для задания 2}{fig:code_task_2}
  \begin{minipage}{.75\textwidth}
    \lstinputlisting[language=Python,basicstyle=\fontsize{10}{10}\linespread{1}\selectfont\ttfamily]{code/task2.py}
  \end{minipage}
\end{vvsu_figure}

% Подглава - Задание 3
\subsection{Задание 3}

Программа ищет в списке все вхождения числа 3 и заменяет их на 30. Цикл проходит по индексам списка, сравнивая каждый элемент с числом 3. Если совпадение найдено, значение по этому индексу заменяется на 30. На выходе получается модифицированный список, в котором все тройки заменены на заданное число. На рисунке \ref{fig:code_task_3} представлен код программы.

\begin{vvsu_figure}{Листинг программы для задания 3}{fig:code_task_3}
  \begin{minipage}{.75\textwidth}
    \lstinputlisting[language=Python,basicstyle=\fontsize{10}{10}\linespread{1}\selectfont\ttfamily]{code/task3.py}
  \end{minipage}
\end{vvsu_figure}

% Подглава - Задание 4
\subsection{Задание 4}

Данный фрагмент кода пытается отсортировать кортеж. Он помещает операцию сортировки в блок try, который преобразует кортеж в отсортированный список с помощью функции sorted, а затем снова создаёт из него кортеж. Если в процессе возникает какое-либо исключение (например, если бы элементы кортежа были несравнимы), программа переходит в блок except, и переменной result просто присваивается исходный кортеж без изменений. В конце выводится полученный результат.На рисунке \ref{fig:code_task_4} представлен код решения.

\begin{vvsu_figure}{Листинг программы для задания 4}{fig:code_task_4}
  \begin{minipage}{.75\textwidth}
    \lstinputlisting[language=Python,basicstyle=\fontsize{10}{10}\linespread{1}\selectfont\ttfamily]{code/task4.py}
  \end{minipage}
\end{vvsu_figure}


% Подглава - Задание 5
\subsection{Задание 5}

Код находит товары с минимальной и максимальной ценой в словаре, где ключи — названия продуктов, а значения — их цены. Для этого используются функции min и max с параметром key, который указывает, что сравнивать нужно по значениям словаря (ценам). В результате определяются ключи (названия товаров), соответствующие минимальной и максимальной цене, после чего выводятся оба значения вместе с ценами. На рисунке \ref{fig:code_task_5} представлен код программы.

\begin{vvsu_figure}{Листинг программы для задания 5}{fig:code_task_5}
  \begin{minipage}{.75\textwidth}
    \lstinputlisting[language=Python,basicstyle=\fontsize{10}{10}\linespread{1}\selectfont\ttfamily]{code/task5.py}
  \end{minipage}
\end{vvsu_figure}

% Подглава - Задание 6
\subsection{Задание 6}

Программа создаёт словарь из списка разнотипных элементов, используя их строковые представления в качестве ключей. Для каждого элемента в исходном списке генерируется пара "ключ-значение", где ключом становится результат преобразования элемента в строку с помощью функции str, а значением — сам исходный элемент. Таким образом, в итоговом словаре ключи — это строки, а значения сохраняют оригинальные типы данных. На рисунке \ref{fig:code_task_6} представлен код программы.

\begin{vvsu_figure}{Листинг программы для задания 6}{fig:code_task_6}
  \begin{minipage}{.75\textwidth}
    \lstinputlisting[language=Python,basicstyle=\fontsize{10}{10}\linespread{1}\selectfont\ttfamily]{code/task6.py}
  \end{minipage}
\end{vvsu_figure}

% Подглава - Задание 7
\subsection{Задание 7}

Эта программа создаёт англо-русский словарь, а затем инвертирует его, получая русско-английский. Для инверсии используется генератор словаря, который меняет местами ключи и значения исходного словаря. Пользователю предлагается ввести слово на русском языке, и если оно присутствует среди ключей нового словаря, выводится его английский перевод. В противном случае программа сообщает, что слово не найдено. На рисунке \ref{fig:code_task_7} представлен код программы.

\begin{vvsu_figure}{Листинг программы для задания 7}{fig:code_task_7}
  \begin{minipage}{.75\textwidth}
    \lstinputlisting[language=Python,basicstyle=\fontsize{10}{10}\linespread{1}\selectfont\ttfamily]{code/task7.py}
  \end{minipage}
\end{vvsu_figure}

% Подглава - Задание 8
\subsection{Задание 8}

Это реализация расширенной версии игры "Камень-ножницы-бумага" с добавлением ящерицы и Спока. Программа определяет победителя по заданным правилам, которые хранятся в словаре, где ключ — выбор игрока, а значение — список вариантов, которые этот выбор побеждает. После получения выбора пользователя и случайного выбора компьютера программа проверяет условия: ничья, победа пользователя (если выбор компьютера содержится в списке побеждаемых вариантов для выбора пользователя) или победа компьютера. На рисунке \ref{fig:code_task_8} представлен код программы.

\begin{vvsu_figure}{Листинг программы для задания 8}{fig:code_task_8}
  \begin{minipage}{.75\textwidth}
    \lstinputlisting[language=Python,basicstyle=\fontsize{10}{10}\linespread{1}\selectfont\ttfamily]{code/task8.py}
  \end{minipage}
\end{vvsu_figure}

% Подглава - Задание 9
\subsection{Задание 9}
 Код группирует слова из списка фруктов по их первой букве, создавая словарь. Для каждого фрукта определяется его первая буква. Если этой буквы ещё нет среди ключей словаря, создаётся новый ключ с пустым списком в качестве значения. Затем фрукт добавляется в список, соответствующий своей начальной букве. В результате получается словарь, где ключи — это первые буквы, а значения — списки фруктов, которые с них начинаются.На рисунке \ref{fig:code_task_9} представлен код программы.

\begin{vvsu_figure}{Листинг программы для задания 9}{fig:code_task_9}
  \begin{minipage}{.75\textwidth}
    \lstinputlisting[language=Python,basicstyle=\fontsize{10}{10}\linespread{1}\selectfont\ttfamily]{code/task9.py}
  \end{minipage}
\end{vvsu_figure}

% Подглава - Задание 10
\subsection{Задание 10}

 Код вычисляет средний балл для каждого студента и определяет лучшего. На основе списка кортежей (имя, оценки) строится словарь, где ключ — имя студента, а значение — его средний балл. Для нахождения студента с максимальным средним используется функция max с параметром key, которая указывает сравнивать не ключи словаря, а их значения. В результате выводится имя лучшего студента и его средний балл. На рисунке \ref{fig:code_task_10} представлен код программы.

\begin{vvsu_figure}{Листинг программы для задания 10}{fig:code_task_10}
  \begin{minipage}{.75\textwidth}
    \lstinputlisting[language=Python,basicstyle=\fontsize{10}{10}\linespread{1}\selectfont\ttfamily]{code/task10.py}
  \end{minipage}
\end{vvsu_figure}

Спасибо за внимание !

\end{document}